\subsection{同态映射的代数独立性}

定理2. --- 设K为无限域，L为K的扩域，A为K上的代数。设 \(u_1, \ldots, u_n\) 是A到L的不同K-代数同态，且 f 是多项式环 L{[}X₁, \ldots，Xₙ{]} 中的多项式。如果我们有 \(f(u_1(x), \ldots, u_n(x)) = 0\) （对所有 \(x \in \mathbf{A}\)），则 \(f = 0\) 。

设 B 为 \(L^{n}\) 中形如 \((u_{1}(x), \ldots, u_{n}(x))\) （其中 \(x \in A\)）的元素之集。根据定理 1，不存在元素不全为零的序列 \((\alpha_{1}, \ldots, \alpha_{n})\) 使得 \(\sum_{i=1}^{n} \alpha_{i} u_{i}(x) = 0\) 对所有 \(x \in A\) 成立；因此（II，第301页，定理7）B 生成

向量空间 \(\mathbf{L}^n\) （在 \(\mathbf{L}\) 上）。因此存在元素 \(a_1, \ldots, a_n\) （属于 \(\mathbf{A}\)）使得矩阵 \((u_i(a_j))_{1 \leqslant i, j \leqslant n}\) 是可逆的。

我们如下定义多项式 \(g \in L[Y_1, \ldots, Y_n]\) ：

\[
g \left(\mathrm{Y} _ {1}, \dots , \mathrm{Y} _ {n}\right) = f \left(\sum_ {j = 1} ^ {n} u _ {1} \left(a _ {j}\right) \mathrm{Y} _ {j}, \dots , \sum_ {j = 1} ^ {n} u _ {n} \left(a _ {j}\right) \mathrm{Y} _ {j}\right).\tag{3}
\]

设 \(y_{1}, \ldots, y_{n}\) 在 K 中；写 \(x = \sum_{i=1}^{n} y_{i} a_{i}\) ，我们有

\[
g \left(y _ {1}, \dots , y _ {n}\right) = f \left(u _ {1} (x), \dots , u _ {n} (x)\right), \quad \text { whence } \quad g \left(y _ {1}, \dots , y _ {n}\right) = 0
\]

根据关于 \(f\) 的假设。由于域 \(K\) 是无限的，我们有 \(g = 0\) （IV，第～18页，推论 2）；现在矩阵 \((u_i(a_j))\) 有一个逆 \((b_{ij})\) ，并且我们有

\[
f \left(\mathrm{X} _ {1}, \dots , \mathrm{X} _ {n}\right) = g \left(\sum_ {j = 1} ^ {n} b _ {1 j} \mathrm{X} _ {j}, \dots , \sum_ {j = 1} ^ {n} b _ {n j} \mathrm{X} _ {j}\right),\tag{4}
\]

由此 \(f = 0\) 。

定理 2 对有限域没有类似结论。于是设 K 为具有 q 个元素的有限域，A = L = K，且 \(f(X) = X^{q} - X\) 。我们有 \(x^{q} = x\) （对所有 \(x \in K\)）（V，第～93页，命题 2）；因此若 u 是 K 的恒等自同构，我们有 \(f(u(x)) = 0\) （对所有 \(x \in K\)），即使 f 不为零。

